% !TeX root=../main.tex

\chapter{دیباچه}
\section{هدف پژوهش}

استفاده از هوش مصنوعی در پاسخگویی به سوالات پزشکی به عنوان یکی از حوزه‌های نوظهور و حیاتی در فناوری و بهداشت، در سال‌های اخیر توجه زیادی را به خود جلب کرده است. در حالی که تحقیقات گسترده‌ای در این زمینه به زبان انگلیسی منتشر شده، در زمان آغاز این پایان‌نامه، حوزه زبان فارسی از پژوهش‌های کافی برخوردار نبود. هدف این پژوهش جمع‌آوری داده‌های پزشکی فارسی، توسعه مدل‌های زبانی کوچک پزشکی به زبان فارسی و ایجاد شیوه‌های محک‌زنی
\footnote{\lr{benchmarks}}
جدید برای سنجش دانش این مدل‌ها است. همچنین، این پژوهش بر توسعه یک مدل زبانی کوچک پزشکی فارسی مبتنی بر استدلال
\footnote{\lr{reasoning}}
تمرکز دارد. مزیت مدل‌های کوچک زبانی نسبت به مدل‌های بزرگ، قابلیت اجرای آن‌ها بر روی دستگاه‌های محلی است. این موضوع اهمیت ویژه‌ای دارد زیرا داده‌های پزشکی معمولا حساس و خصوصی هستند و ارسال آن‌ها به سرورهای خارجی می‌تواند خطرات جدی برای حریم خصوصی بیماران ایجاد کند
\cite{b1}.
\section{کاربرد پژوهش}
\subsection{کاربرد مدل های زبانی پزشکی}
مدل‌های زبانی در سال‌های اخیر با استفاده از داده‌های بسیار گسترده‌تر و معماری‌های پیشرفته‌تر به پیشرفت‌های چشمگیری دست یافته‌اند. این مدل‌ها توانایی درک بهتر مفاهیم، تولید متن‌های طبیعی‌تر و پاسخ‌دهی دقیق‌تر به سوالات را پیدا کرده‌اند. 

 این پیشرفت‌ها منجر به افزایش چشمگیر کاربرد هوش مصنوعی
\footnote{\lr{artificial intelligence}}
در حوزه‌های مختلف، به‌ویژه در زمینه پزشکی، شده است.
\cite{b2}
امروزه در حوزه پزشکی، مدل‌های زبانی مبتنی بر یادگیری ژرف
\footnote{\lr{deep learning}}
نقش مهمی در تحلیل داده‌های پزشکی، بهبود دقت تشخیص بیماری‌ها، ارائه پیشنهادهای درمانی دقیق‌تر و افزایش کیفیت مراقبت از بیماران ایفا می‌کنند. علاوه بر این، این فناوری، به بهینه‌سازی سیستم‌های اداری و کاهش بار کاری کادر درمانی کمک شایانی کرده است. به عنوان مثال، مدل‌های هوش مصنوعی قادرند با تحلیل داده‌های حاصل از پرونده‌های پزشکی، الگوهای مرتبط با بیماری‌ها را شناسایی کنند و اطلاعات ارزشمندی را برای تصمیم‌گیری سریع‌تر و دقیق‌تر در اختیار پزشکان قرار دهند. 
این مدل‌ها همچنین می‌توانند نقش مهمی در تکمیل مشاوره‌های پزشکی ایفا کرده و به پزشکان در ارائه اطلاعات دقیق‌تر و سریع‌تر کمک کنند.و حتی شاید در آینده ای نه چندان دور بتوانند جای پزشکان را در مشاوره های پزشکی بگیرند.
\cite{b3}

این تحول نه تنها به افزایش کارایی و بهره‌وری در سیستم‌های درمانی منجر شده است، بلکه تجربه کلی بیماران را نیز بهبود بخشیده و امکان ارائه خدمات درمانی بهتر و موثرتر را فراهم کرده است. به همین دلیل، توسعه و استفاده از مدل های زبانی پزشکی
\footnote{\lr{medical language models}}
، همچنان مورد توجه پژوهشگران و متخصصان قرار دارد.

\subsection{کاربرد مدل های زبانی پزشکی دارای قابلیت استدلال}
پزشکی به‌عنوان یک علم تجربی و کاربردی، بیش از هر چیز به استدلال دقیق وابسته است. تشخیص بیماری، انتخاب درمان، پیش‌بینی پیامد و حتی آموزش پزشکی، همگی بر پایهٔ زنجیره‌های منطقی، تحلیل داده‌های بالینی و استنتاج از شواهد بنا شده‌اند. از آن روی که پزشکی علم وابسته به استدلال است، توسعه یک مدل زبانی که دارای قابلیت‌های استدلال برتری باشد می‌تواند بسیار تحول‌آفرین باشد. این مدل‌ها نه‌تنها می‌توانند اطلاعات را بازیابی کنند، بلکه می‌توانند دلیل بیاورند، احتمالات را مقایسه کنند، تناقضات را شناسایی کنند و تصمیمات را توجیه‌پذیر سازند. بنابراین توسعه این نوع از مدل های زبانی در حوزه پزشکی بسیار مفید است.
\cite{b4}
\subsection{کاربرد مدل های زبانی پزشکی فارسی}
علیرغم پیشرفت‌های چشمگیر در توسعه مدل‌های زبانی پزشکی به زبان انگلیسی، در حوزه زبان فارسی هنوز کار چندانی صورت نگرفته است. این در حالی است که در سرتاسر جهان میلیون ها نفر تنها قادر به استفاده از این زبان هستند؛
بنابراین تلاش برای توسعه یک مدل زبانی پزشکی در زبان فارسی میتواند گامی رو به جلو در ارتباطات و خدمات درمانی کشور های فارسی زبان باشد. 
\section{مراحل انجام پایان نامه}
همانطور که در جدول 
\ref{table:phases_information}
این پایان‌نامه در دو فاز اصلی طراحی و اجرا شده است. فاز نخست به جمع‌آوری دادگان پزشکی فارسی و توسعه مدلی با نام «گائوکرنا وی» اختصاص دارد که فاقد توانایی استدلال بوده و بیشتر بر درک سیستم یک 
\footnote{
در علم رفتارشناسی به درک سریع، شهودی و بدون نیاز به تفکر ژرف درک سیستم یک و به درک آهسته، غیر شهودی و نیازمند استدلال درک سیستم دو میگویند.در فصل بعدی به این موضوع بیشتر پرداخته خواهد شد.
}
زبان تمرکز دارد. از این فاز، مقاله‌ای با عنوان "اهرم قرار دادن داده‌های آنلاین برای بهبود دانش پزشکی یک مدل زبانی کوچک پزشکی فارسی" 
\footnote{\lr{Leveraging Online Data to Enhance Medical Knowledge in a Small Persian Language Model}}
استخراج شده است که به تشریح فرآیند جمع‌آوری داده‌ها و نحوه بهینه‌سازی دانش پزشکی مدل می‌پردازد. در فاز دوم این پژوهش ابتدا تکنیک های جدیدی برای ارتقای توانایی استدلال و درک سیستم دو مدل معرفی شده و سپس مدل «گائوکرنا آر» در این فاز توسعه داده شده است.
از این فاز نیز، مقاله‌ای با عنوان "بهبود مهارت‌های استدلال در مدل‌های زبانی پزشکی فارسی کوچک می‌تواند از آموزش داده‌های بزرگ‌مقیاس پیشی بگیرد"
\footnote{\lr{Enhancing Reasoning Skills in Small Persian Medical Language Models Can Outperform Large-Scale Data Training}}
استخراج شده است.

مقاله نخست در سی و دومین کنفرانس ملی و دهمین کنفرانس بین‌المللی مهندسی پزشکی ایران و مقاله دوم در یازدهمین کنفرانس بین‌المللی پردازش سیگنال و سیستم‌های هوشمند منتشر شده است.
 
\begin{table}[h!]
	\centering
	\begin{tabular}{|c|c|c|}
		\hline
		\textbf{\lr{gaokerena-R}}& \textbf{\lr{gaokerena-V}} &   \\ \hline
		\lr{mehrdadghassabi/gaokerena-R}   &     \lr{mehrdadghassabi/gaokerena-V}  & مخزن گیت هاب      \\ \hline
		\lr{gaokerena/gaokerena-r1.0}    &     \lr{gaokerena/gaokerena-v1.0}  & مخزن پارامتر ها      \\ \hline
		\lr{arxiv.org/pdf/2510.20059}  &    \lr{arxiv.org/pdf/2505.16000}   & پیوند مقاله      \\ \hline
		70 دلار  &      300 دلار   & هزینه     \\ \hline
		بله  &     خیر   & توانایی استدلال     \\ \hline
		یازدهمین کنفرانس بین‌المللی &   سی و دومین کنفرانس ملی   & منتشر شده در    \\ 
		 پردازش سیگنال و سیستم‌های هوشمند &    و دهمین کنفرانس بین‌المللی   &    \\ 
		&   مهندسی پزشکی ایران   &    \\ \hline
	\end{tabular}
	\caption{اطلاعات دو فاز پایان نامه}
	\label{table:phases_information}
\end{table}


\section{ساختار پایان نامه}
در این پایان‌نامه، ساختار فصل‌ها به‌گونه‌ای طراحی شده است که مراحل مختلف پژوهش به‌صورت منظم و هدفمند ارائه شوند. فصل دوم به بررسی ادبیات موضوع مورد تحقیق اختصاص دارد، در حالی که فصل سوم به مرور کارهای پیشین پرداخته و مطالعات انجام‌شده در زمینه‌های مرتبط را بررسی می‌کند.
فصل چهارم به دلیل کمبود داده‌های پزشکی در حوزه زبان فارسی، فرآیند گردآوری و آماده‌سازی این داده‌ها را به‌طور دقیق تشریح خواهد کرد. سپس در فصل پنجم، با استفاده از داده‌های معرفی‌شده در فصل سوم، مدل اولیه‌ای با نام «گائوکرنا وی»
\footnote{نام گائوکرنا از درختی افسانه‌ای الهام گرفته شده است که در روایات اساطیری زرتشتی به‌عنوان نماد شفادهی و جاودانگی شناخته می‌شود.}
معرفی و تحلیل می‌شود.
در ادامه، فصل ششم به معرفی تکنیک‌هایی برای بهبود توانایی استدلال یک مدل زبانی می‌پردازد و مدل پیشرفته‌تری به نام «گائوکرنا آر» را معرفی کرده و ویژگی‌های آن را به‌تفصیل شرح می‌دهد. در نهایت، فصل پایانی به جمع‌بندی نتایج پژوهش و ارائه پیشنهاداتی برای تحقیقات آینده اختصاص دارد.

